\subsubsection{Gauss-Jordan (mod p)}
\begin{lstlisting}[style=mystyle, language=C++]
// TC: O(N^3) where N is the matrix dimension
// usa: resolver sistemas lineales, calcular determinante y rango
#include <bits/stdc++.h>
using namespace std;

const long long MOD = 998244353LL;

// returns rank; mutates a in-place to row-echelon form
int gaussJordan(vector<vector<long long>>& a){
	int n = a.size(), m = a[0].size(), rank = 0;
	int row = 0;
	for(int col=0; col<m && row<n; col++){
		int sel = row;
		for(int i=row;i<n;i++) if(a[i][col]){ sel = i; break; }
		if(!a[sel][col]) continue;
		swap(a[sel], a[row]);
		long long inv = modpow(a[row][col], MOD - 2, MOD);
		for(int j=col;j<m;j++) a[row][j] = a[row][j] * inv % MOD;
		for(int i=0;i<n;i++){
			if(i != row && a[i][col]){
				long long factor = a[i][col];
				for(int j=col;j<m;j++){
					a[i][j] = (a[i][j] - factor * a[row][j] % MOD + MOD) % MOD;
				}
			}
		}
		row++; rank++;
	}
	return rank;
}

long long determinant(vector<vector<long long>> a){
	int n = a.size();
	long long det = 1;
	for(int i=0;i<n;i++){
		int sel = i;
		for(int j=i;j<n;j++) if(a[j][i]){ sel = j; break; }
		if(!a[sel][i]) return 0;
		if(sel != i){ swap(a[sel], a[i]); det = (MOD - det) % MOD; }
		det = det * a[i][i] % MOD;
		long long inv = modpow(a[i][i], MOD - 2, MOD);
		for(int j=i+1;j<n;j++){
			long long factor = a[j][i] * inv % MOD;
			for(int k=i;k<n;k++){
				a[j][k] = (a[j][k] - factor * a[i][k] % MOD + MOD) % MOD;
			}
		}
	}
	return det;
}

// --- Resolver Ax = b usando gaussJordan sobre la matriz aumentada ---
// gaussJordan devuelve el RANGO r de la matriz aumentada [A|b].
//   - Si tras la REF existe alguna fila i con todos los coeficientes
//     en cero pero el termino independiente aug[i][n] != 0, el
//     sistema es INCONSISTENTE: no hay solucion.
//   - Si r == n (rango == #incognitas), hay SOLUCION UNICA y vive
//     en la ultima columna de la RREF: x[i] = aug[i][n].
//   - Si r < n, quedan (n - r) variables libres: hay INFINITAS
//     soluciones. La ultima columna da una solucion particular
//     poniendo las libres en 0.
// Retorna:
//   0 -> sin solucion
//   1 -> solucion unica (escrita en 'sol')
//   2 -> infinitas soluciones (escribe una particular en 'sol')
int solveSystem(vector<vector<long long>> A, vector<long long> b,
				vector<long long>& sol){
	int n = (int)A.size();
	// 1) construir la matriz aumentada [A | b] de tamano n x (n+1)
	vector<vector<long long>> aug(n, vector<long long>(n + 1));
	for(int i = 0; i < n; i++){
		for(int j = 0; j < n; j++) aug[i][j] = A[i][j];
		aug[i][n] = b[i];
	}

	// 2) RREF + rango
	int r = gaussJordan(aug);

	// 3) checar inconsistencia: fila nula con termino independiente != 0
	for(int i = r; i < n; i++)
		if(aug[i][n] != 0) return 0;   // sin solucion

	// 4) rango < #incognitas -> variables libres -> infinitas soluciones
	if(r < n) return 2;

	// 5) solucion unica: se lee directo de la ultima columna
	sol.assign(n, 0);
	for(int i = 0; i < n; i++) sol[i] = aug[i][n];
	return 1;
}

int main(){
	vector<vector<long long>> A = {{1,2,3},{4,5,6},{7,8,10}};
	int r = gaussJordan(A);
	cout << "rank = " << r << "\n";
	cout << "det = " << determinant({{1,2},{3,4}}) << "\n";  // -2 mod p

	// Ejemplo: 2x + y = 5 ,  x - y = 1   ->  x = 2, y = 1
	vector<vector<long long>> M = {{2,1},{1,-1}};
	vector<long long> b = {5,1}, x;
	int status = solveSystem(M, b, x);   // status = 1, x = {2, 1}
	for(long long v : x) cout << v << " ";
	cout << "\n";
	return 0;
}
\end{lstlisting}
